
\documentclass{article}
\usepackage{colortbl}
\usepackage{makecell}
\usepackage{multirow}
\usepackage{supertabular}

\begin{document}

\newcounter{utterance}

\centering \large Interaction Transcript for game `ifeval', experiment `ifeval', episode 54 with refineonly{-}Qwen3.5{-}27B.
\vspace{24pt}

{ \footnotesize  \setcounter{utterance}{1}
\setlength{\tabcolsep}{0pt}
\begin{supertabular}{c@{$\;$}|p{.15\linewidth}@{}p{.15\linewidth}p{.15\linewidth}p{.15\linewidth}p{.15\linewidth}p{.15\linewidth}}
    \# & \multicolumn{2}{c}{Player} && \multicolumn{2}{c}{Game Master} \\
    \hline

    \theutterance \stepcounter{utterance}  
    & & & \multicolumn{4}{p{0.6\linewidth}}{
        \cellcolor[rgb]{0.9,0.9,0.9}{
            \makecell[{{p{\linewidth}}}]{
                \texttt{\tiny{[P1$\langle$GM]}}
                \texttt{Take the text below as a starting point, and make it a complete article: "You may have to meet with a helper to work out a parenting plan. The first would be to estimate how many time you have everyday for parenting, and is that enough...."} \\
\texttt{Avoid using the following keywords: sleep, cook, feed} \\
\texttt{Mention the keyword 'schedule' for more than 5 times.} \\
            }
        }
    }
     \\ \\

    \theutterance \stepcounter{utterance}  
    & & \multicolumn{4}{p{0.6\linewidth}}{
        \cellcolor[rgb]{0.95,0.95,0.95}{
            \makecell[{{p{\linewidth}}}]{
                \texttt{\tiny{[P1$\leftrightarrow$P]}}
                \texttt{{[}DRAFT{]}} \\
\texttt{\# Crafting a Sustainable Parenting Plan: A Guide to Time and Balance} \\
\\ 
\texttt{You may have to meet with a helper to work out a parenting plan. The first step in this crucial process is to estimate how much time you have every day for parenting and determine if that amount is sufficient for your family's needs. This initial assessment sets the foundation for a realistic and effective **schedule** that accommodates both your professional obligations and your role as a caregiver.} \\
\\ 
\texttt{When you sit down with a mediator or counselor, the primary goal is to create a **schedule** that reflects the unique rhythm of your household. It is not enough to simply divide hours; you must analyze the quality of those hours. A rigid **schedule** often fails because it does not account for unexpected events, such as a child's illness or a sudden work meeting. Therefore, flexibility must be woven into the very fabric of your **schedule**.} \\
\\ 
\texttt{During your meeting, you will likely discuss how to allocate time for educational activities, play, and emotional connection. You need to ask yourself if your current **schedule** allows for meaningful interaction or if it is merely a series of logistical tasks. If the **schedule** feels too packed, you may need to delegate certain responsibilities or adjust your work hours to ensure you are fully present.} \\
\\ 
\texttt{It is also vital to consider the long{-}term viability of the plan. A **schedule** that works for a week might become unsustainable over a month. You should review your **schedule** regularly with your helper to make necessary adjustments as your children grow and their needs evolve. By prioritizing a balanced **schedule**, you ensure that your parenting time is not just about being physically present, but about being emotionally available.} \\
\\ 
\texttt{Ultimately, the success of your parenting plan relies on a **schedule** that everyone in the family can understand and follow. Whether you are navigating a separation or simply reorganizing your family life, a well{-}structured **schedule** provides the stability children need to thrive. Remember, the goal is not perfection, but a consistent and loving routine that supports your family's well{-}being.} \\
            }
        }
    }
    & \\ \\

    \theutterance \stepcounter{utterance}  
    & \multicolumn{4}{p{0.6\linewidth}}{
        \cellcolor[rgb]{0.561,0.737,0.561}{
            \makecell[{{p{\linewidth}}}]{
                \texttt{\tiny{[P1$\rangle$GM]}}
                \texttt{\# Crafting a Sustainable Parenting Plan: A Guide to Time and Balance} \\
\\ 
\texttt{You may have to meet with a helper to work out a parenting plan. The first step in this crucial process is to estimate how much time you have every day for parenting and determine if that amount is sufficient for your family's needs. This initial assessment sets the foundation for a realistic and effective **schedule** that accommodates both your professional obligations and your role as a caregiver.} \\
\\ 
\texttt{When you sit down with a mediator or counselor, the primary goal is to create a **schedule** that reflects the unique rhythm of your household. It is not enough to simply divide hours; you must analyze the quality of those hours. A rigid **schedule** often fails because it does not account for unexpected events, such as a child's illness or a sudden work meeting. Therefore, flexibility must be woven into the very fabric of your **schedule**.} \\
\\ 
\texttt{During your meeting, you will likely discuss how to allocate time for educational activities, play, and emotional connection. You need to ask yourself if your current **schedule** allows for meaningful interaction or if it is merely a series of logistical tasks. If the **schedule** feels too packed, you may need to delegate certain responsibilities or adjust your work hours to ensure you are fully present.} \\
\\ 
\texttt{It is also vital to consider the long{-}term viability of the plan. A **schedule** that works for a week might become unsustainable over a month. You should review your **schedule** regularly with your helper to make necessary adjustments as your children grow and their needs evolve. By prioritizing a balanced **schedule**, you ensure that your parenting time is not just about being physically present, but about being emotionally available.} \\
\\ 
\texttt{Ultimately, the success of your parenting plan relies on a **schedule** that everyone in the family can understand and follow. Whether you are navigating a separation or simply reorganizing your family life, a well{-}structured **schedule** provides the stability children need to thrive. Remember, the goal is not perfection, but a consistent and loving routine that supports your family's well{-}being.} \\
            }
        }
    }
    & & \\ \\

    \theutterance \stepcounter{utterance}  
    & & & \multicolumn{2}{p{0.3\linewidth}}{
        \cellcolor[rgb]{0.9,0.9,0.9}{
            \makecell[{{p{\linewidth}}}]{
                \texttt{\tiny{[GM$|$GM]}}
                {'keywords:forbidden_words': {'forbidden_words': ['sleep', 'cook', 'feed']}, 'keywords:frequency': {'relation': 'at least', 'keyword': 'schedule', 'frequency': 6}}
            }
        }
    }
    & & \\ \\

    \theutterance \stepcounter{utterance}  
    & & & \multicolumn{2}{p{0.3\linewidth}}{
        \cellcolor[rgb]{0.9,0.9,0.9}{
            \makecell[{{p{\linewidth}}}]{
                \texttt{\tiny{[GM$|$GM]}}
                \texttt{game\_result = WIN} \\
            }
        }
    }
    & & \\ \\

\end{supertabular}
}

\end{document}
